\scanpage{201}\mspage{146.}
et \unsure{compte} tenu que le composé de la première ligne horizontale est $i_{X'}$, par construction de $\alpha_{X'}$, on est ramené à la commutativité de
\diagram{\begin{tikzcd}[ampersand replacement=\&,column sep=large]
 \& X\times_S S'=X' \arrow[r,no head,"i_{X'}"] \arrow[d,"{\unsure{j|i}_X\times_S S'}"'] \& \Phi_S(X') \arrow[d,"{\Phi_S(\mathrm{pr}_1^X)}"]\\
 (*) \& \Phi_S(X)\times_S S' \arrow[r,"{\mathrm{pr}_1^{\Phi_S(X)}}"'] \& \Phi_S(X).
\end{tikzcd}}
Or le morphisme $i_X:X\longrightarrow\Phi_S(X)$ dans $\mathrm{Cat}_{/S}$ donne lieu au carré commutatif (comme \unsure{toute} flèche $i:X\longrightarrow Y$)
\diagram{\begin{tikzcd}[ampersand replacement=\&,column sep=large]
 X\times_S S' \arrow[r,"\mathrm{pr}_1^X"] \arrow[d,"{i_X\times_S S'}"'] \& X \arrow[d,"\unsure{i'_X|i_X}"]\\
 \Phi_S(X)\times_S S' \arrow[r,"{\mathrm{pr}_1^{\Phi_S(X)}}"] \& \Phi_S(X),
\end{tikzcd}}
de sorte que le composé inférieur de $(*)$ n'est autre que celui de $(**)$, donc égal au composé sup. de ce dernier, et la comm. de $(*)$ revient à celle de
\diagram{\begin{tikzcd}[ampersand replacement=\&,column sep=large]
 X\times_S S'=X' \arrow[r,"i_{X'}"] \arrow[d,"\mathrm{pr}_1^X"'] \& \Phi_S(X') \arrow[d,"{\Phi_S(\mathrm{pr}_1^X)}"]\\
 X \arrow[r,"i_X"] \& \Phi_S(X),
\end{tikzcd}}
qui est évidente par fonctorialité de \inserted{$\underbrace{(Z\xrightarrow{i_Z}\Phi_S(Z))}_{\deleted{\text{\unsure{l'hom}}}\,\text{\unsure{en }}Z}$}, appliquée à la flèche $\mathrm{pr}_1^X:X'\longrightarrow X$ de $\mathrm{Cat}_{/S}$.
% archive sheet number: 200
